% ═══════════════════════════════════════════════════════════════
% ARBEITSBLATT: Das Ohmsche Gesetz (Klasse 8)
% Differenziert: * (Basis) / ** (Standard) / *** (Erweitert)
% ═══════════════════════════════════════════════════════════════

\documentclass[11pt, a4paper]{article}

% ═══════════════════════════════════════════════════════════════
% PAKETE
% ═══════════════════════════════════════════════════════════════

\usepackage[utf8]{inputenc}
\usepackage[T1]{fontenc}
\usepackage[ngerman]{babel}
\usepackage{helvet}
\renewcommand{\familydefault}{\sfdefault}

\usepackage[margin=2cm, top=2.5cm, bottom=2cm]{geometry}
\usepackage{amsmath, amssymb}
\usepackage{siunitx}
\sisetup{locale = DE, per-mode = symbol}

\usepackage{graphicx}
\usepackage{tikz}
\usetikzlibrary{circuits.ee.IEC, positioning}
\usepackage{enumitem}
\usepackage{tabularx}
\usepackage{booktabs}
\usepackage{array}
\usepackage{multicol}

\usepackage{xcolor}
\usepackage{tcolorbox}
\usepackage{fancyhdr}
\usepackage{lastpage}

% Kompaktere Listen
\setlist{nosep, leftmargin=1.5em}

% ═══════════════════════════════════════════════════════════════
% FARBEN
% ═══════════════════════════════════════════════════════════════

\definecolor{physik}{HTML}{2c5aa0}
\definecolor{grau}{HTML}{888888}
\definecolor{hellgrau}{HTML}{f0f0f0}
\definecolor{success}{HTML}{2a7a4b}
\definecolor{stern1}{HTML}{2a7a4b}    % Grün für *
\definecolor{stern2}{HTML}{b35c00}    % Orange für **
\definecolor{stern3}{HTML}{c41e3a}    % Rot für ***

% ═══════════════════════════════════════════════════════════════
% VARIABLEN
% ═══════════════════════════════════════════════════════════════

\newcommand{\thema}{Das Ohmsche Gesetz}
\newcommand{\fach}{Physik}
\newcommand{\klasse}{8}

% ═══════════════════════════════════════════════════════════════
% KOPF- UND FUSSZEILE
% ═══════════════════════════════════════════════════════════════

\pagestyle{fancy}
\fancyhf{}
\renewcommand{\headrulewidth}{0.4pt}
\renewcommand{\footrulewidth}{0pt}

\lhead{\small\textcolor{physik}{\textbf{\fach}} Kl.\,\klasse{} -- \thema}
\rhead{\small Name: \rule{5cm}{0.4pt}}

\cfoot{\small\thepage/\pageref{LastPage}}

% ═══════════════════════════════════════════════════════════════
% BEFEHLE
% ═══════════════════════════════════════════════════════════════

% Lücke für Eintragungen
\newcommand{\luecke}[1]{\rule{#1}{0.4pt}}

% Antwortzeile
\newcommand{\antwort}[1]{\vspace{#1}\hrule\vspace{2pt}}

% Formelbox
\newtcolorbox{formelbox}{
    colback=hellgrau,
    colframe=physik,
    boxrule=1pt,
    arc=0pt,
    left=8pt, right=8pt, top=8pt, bottom=8pt
}

% Merksatz
\newtcolorbox{merksatz}{
    colback=hellgrau,
    colframe=success,
    boxrule=1pt,
    arc=0pt,
    left=8pt, right=8pt, top=6pt, bottom=6pt,
    fonttitle=\bfseries,
    title=Merksatz
}

% Niveau-Marker
\newcommand{\niveauEins}{\textcolor{stern1}{\textbf{(*)}}}
\newcommand{\niveauZwei}{\textcolor{stern2}{\textbf{(**)}}}
\newcommand{\niveauDrei}{\textcolor{stern3}{\textbf{(***)}}}

% Punkteangabe
\newcommand{\punkte}[1]{\hfill\textcolor{grau}{\small[#1 BE]}}

% ═══════════════════════════════════════════════════════════════
% DOKUMENT
% ═══════════════════════════════════════════════════════════════

\begin{document}

% ═══════════════════════════════════════════════════════════════
% KOPFBEREICH
% ═══════════════════════════════════════════════════════════════

\begin{center}
    {\Large\textbf{\textcolor{physik}{Arbeitsblatt: \thema}}}
\end{center}

\vspace{0.5em}

\begin{tcolorbox}[colback=hellgrau, colframe=grau, boxrule=0.5pt, arc=0pt]
\small
\textbf{Niveaus:}
\niveauEins{} Basis \quad
\niveauZwei{} Standard \quad
\niveauDrei{} Erweitert
\hfill
\textbf{Gesamtpunktzahl:} \luecke{1cm} / 25 BE
\end{tcolorbox}

\vspace{1em}

% ═══════════════════════════════════════════════════════════════
% TEIL 1: MERKSATZ (zum Übertragen aus dem Lernpfad)
% ═══════════════════════════════════════════════════════════════

\section*{1. Formel und Merksatz}

\begin{formelbox}
\textbf{Das Ohmsche Gesetz:}

\vspace{1em}
\begin{center}
{\Large $U = R \cdot I$}
\end{center}

\vspace{0.5em}

\begin{tabular}{lll}
$U$ & = & Spannung in \luecke{2cm} (Einheit: \luecke{1.5cm}) \\[0.8em]
$R$ & = & \luecke{3cm} in Ohm (Einheit: \luecke{1.5cm}) \\[0.8em]
$I$ & = & Stromstärke in \luecke{2cm} (Einheit: \luecke{1.5cm}) \\
\end{tabular}
\end{formelbox}

\vspace{0.5em}

\begin{merksatz}
Je \luecke{2cm} die Spannung, desto \luecke{2cm} die Stromstärke.

Der Widerstand $R$ bleibt dabei \luecke{2.5cm} (bei ohmschen Widerständen).
\end{merksatz}

\vspace{1em}

% ═══════════════════════════════════════════════════════════════
% TEIL 2: MESSWERTE (aus dem Experiment)
% ═══════════════════════════════════════════════════════════════

\section*{2. Messwerte aus dem Experiment \niveauEins}

\textbf{Widerstand:} $R = \SI{20}{\ohm}$

\vspace{0.5em}

\begin{center}
\begin{tabular}{|c|c|c|c|c|c|}
\hline
\textbf{Spannung $U$} & \SI{2}{\volt} & \SI{4}{\volt} & \SI{6}{\volt} & \SI{8}{\volt} & \SI{10}{\volt} \\
\hline
\textbf{Stromstärke $I$} & \luecke{1.2cm} & \luecke{1.2cm} & \luecke{1.2cm} & \luecke{1.2cm} & \luecke{1.2cm} \\
\hline
\end{tabular}
\end{center}

\vspace{0.5em}

\textbf{Zeichne die Messwerte als Punkte ein und verbinde sie mit einer Geraden:}

\vspace{0.5em}

\begin{center}
\begin{tikzpicture}[scale=0.8]
    % Kästchenraster
    \draw[step=0.5cm, gray!30, very thin] (0,0) grid (8,6);
    \draw[step=1cm, gray!50, thin] (0,0) grid (8,6);

    % Achsen
    \draw[thick, ->] (0,0) -- (8.3,0) node[right] {$U$ in V};
    \draw[thick, ->] (0,0) -- (0,6.3) node[above] {$I$ in A};

    % X-Achsen-Beschriftung
    \foreach \x in {0,2,4,6,8,10} {
        \draw (\x/1.25,0) -- (\x/1.25,-0.1);
        \node[below] at (\x/1.25,-0.1) {\small\x};
    }

    % Y-Achsen-Beschriftung
    \foreach \y/\label in {0/0, 1/0.1, 2/0.2, 3/0.3, 4/0.4, 5/0.5} {
        \draw (0,\y) -- (-0.1,\y);
        \node[left] at (-0.1,\y) {\small\label};
    }
\end{tikzpicture}
\end{center}

\vspace{0.5em}

% ═══════════════════════════════════════════════════════════════
% TEIL 3: AUFGABEN NIVEAU *
% ═══════════════════════════════════════════════════════════════

\section*{3. Aufgaben}

\subsection*{Niveau \niveauEins{} -- Basis}

\textbf{Aufgabe 1:} Ergänze den Lückentext. \punkte{3}

Wenn man die Spannung am Widerstand \textbf{verdoppelt}, dann \luecke{3cm} sich auch

die \luecke{3cm}.

Man sagt: Die Stromstärke ist \luecke{3cm} zur Spannung.

\vspace{1.5em}

\textbf{Aufgabe 2:} Berechne die Spannung $U$. \punkte{2}

\begin{tcolorbox}[colback=white, colframe=physik!50, boxrule=0.5pt, arc=0pt]
Gegeben: $R = \SI{10}{\ohm}$, \quad $I = \SI{2}{\ampere}$

\vspace{0.5em}
Gesucht: $U = \,?$

\vspace{0.5em}
Formel: $U = R \cdot I$

\vspace{0.5em}
Einsetzen: $U = $ \luecke{1.5cm} $\cdot$ \luecke{1.5cm} $=$ \luecke{2cm}
\end{tcolorbox}

\vspace{1em}

% ═══════════════════════════════════════════════════════════════
% TEIL 4: AUFGABEN NIVEAU **
% ═══════════════════════════════════════════════════════════════

\subsection*{Niveau \niveauZwei{} -- Standard}

\textbf{Aufgabe 3:} Berechne die fehlenden Größen. \punkte{6}

\begin{enumerate}[label=\alph*)]
    \item $R = \SI{50}{\ohm}$, \quad $I = \SI{0,2}{\ampere}$ \quad $\Rightarrow$ \quad $U = $ \luecke{3cm}

    \vspace{1em}

    \item $U = \SI{12}{\volt}$, \quad $R = \SI{4}{\ohm}$ \quad $\Rightarrow$ \quad $I = $ \luecke{3cm}

    \textit{Hinweis: Forme zuerst um: $I = \dfrac{U}{R}$}

    \vspace{1em}

    \item $U = \SI{10}{\volt}$, \quad $I = \SI{0,5}{\ampere}$ \quad $\Rightarrow$ \quad $R = $ \luecke{3cm}

    \textit{Hinweis: Forme zuerst um: $R = \dfrac{U}{I}$}
\end{enumerate}

\vspace{1em}

\textbf{Aufgabe 4:} Ordne zu: Welche Einheit gehört zu welcher Größe? \punkte{3}

\begin{center}
\begin{tabular}{ccc}
\textbf{Größe} & & \textbf{Einheit} \\[0.5em]
Spannung $U$ & \luecke{2cm} & Ampere (A) \\
Stromstärke $I$ & \luecke{2cm} & Ohm ($\Omega$) \\
Widerstand $R$ & \luecke{2cm} & Volt (V) \\
\end{tabular}
\end{center}

\vspace{1em}

% ═══════════════════════════════════════════════════════════════
% TEIL 5: AUFGABEN NIVEAU ***
% ═══════════════════════════════════════════════════════════════

\subsection*{Niveau \niveauDrei{} -- Erweitert}

\textbf{Aufgabe 5:} Fehlersuche \punkte{4}

Ein Schüler rechnet:

\begin{quote}
``$R = \SI{50}{\ohm}$, $I = \SI{200}{\milli\ampere}$.

Dann ist $U = 50 \cdot 200 = \SI{10000}{\volt}$.''
\end{quote}

a) Was hat der Schüler falsch gemacht?

\antwort{1.5cm}
\antwort{1.5cm}

\vspace{0.5em}

b) Berechne das richtige Ergebnis.

\vspace{2cm}

\vspace{1em}

\textbf{Aufgabe 6:} Transfer \punkte{4}

An einem Widerstand von $R = \SI{20}{\ohm}$ liegt eine Spannung von $U = \SI{6}{\volt}$ an.

\begin{enumerate}[label=\alph*)]
    \item Berechne die Stromstärke $I$.

    \vspace{2cm}

    \item Die Spannung wird auf $U = \SI{12}{\volt}$ verdoppelt.

    Was passiert mit der Stromstärke? Begründe ohne zu rechnen.

    \antwort{1.5cm}
    \antwort{1.5cm}
\end{enumerate}

\vspace{1em}

\textbf{Aufgabe 7:} Begründung \punkte{3}

Erkläre, warum der Widerstand $R$ konstant bleibt, auch wenn man die Spannung ändert.

\textit{Tipp: Was ist der Widerstand -- eine Eigenschaft des Bauteils oder eine Messgröße, die sich ändert?}

\antwort{1.5cm}
\antwort{1.5cm}
\antwort{1.5cm}

\vspace{1.5em}

% ═══════════════════════════════════════════════════════════════
% PUNKTEÜBERSICHT
% ═══════════════════════════════════════════════════════════════

\begin{tcolorbox}[colback=hellgrau, colframe=grau, boxrule=0.5pt, arc=0pt]
\small
\textbf{Punkteübersicht:}

\begin{tabular}{llcl}
\niveauEins{} Basis: & Aufgaben 1--2 & = & 5 BE \\
\niveauZwei{} Standard: & Aufgaben 3--4 & = & 9 BE \\
\niveauDrei{} Erweitert: & Aufgaben 5--7 & = & 11 BE \\
\hline
\textbf{Gesamt:} & & & \textbf{25 BE}
\end{tabular}
\hfill
\textbf{Erreicht:} \luecke{1.5cm} BE
\end{tcolorbox}

% ═══════════════════════════════════════════════════════════════
% FORMELSAMMLUNG (Rückseite oder Hilfsblatt)
% ═══════════════════════════════════════════════════════════════

\newpage

\begin{center}
    {\Large\textbf{\textcolor{physik}{Formelkarte: Ohmsches Gesetz}}}
\end{center}

\vspace{1em}

\begin{formelbox}
\begin{center}
\textbf{Die drei Formen des Ohmschen Gesetzes:}

\vspace{1em}

\begin{tabular}{|c|c|c|}
\hline
\textbf{Spannung berechnen} & \textbf{Stromstärke berechnen} & \textbf{Widerstand berechnen} \\[0.5em]
\hline
& & \\
{\LARGE $U = R \cdot I$} & {\LARGE $I = \dfrac{U}{R}$} & {\LARGE $R = \dfrac{U}{I}$} \\[1em]
\hline
\end{tabular}

\vspace{1.5em}

\textbf{Einheiten:}

\begin{tabular}{lll}
$U$ & in Volt & $[\si{\volt}]$ \\
$I$ & in Ampere & $[\si{\ampere}]$ \\
$R$ & in Ohm & $[\si{\ohm}]$ \\
\end{tabular}

\vspace{1em}

\textbf{Wichtig:} $\SI{1000}{\milli\ampere} = \SI{1}{\ampere}$ \quad (immer in A umrechnen!)

\end{center}
\end{formelbox}

\vspace{2em}

\begin{center}
\textbf{Beispielrechnung:}
\end{center}

\begin{tcolorbox}[colback=white, colframe=physik, boxrule=1pt, arc=0pt, title=\textbf{Beispiel: Stromstärke berechnen}]

\textbf{Gegeben:} $U = \SI{12}{\volt}$, \quad $R = \SI{4}{\ohm}$

\textbf{Gesucht:} $I = \,?$

\textbf{Formel:} $I = \dfrac{U}{R}$

\textbf{Einsetzen:} $I = \dfrac{\SI{12}{\volt}}{\SI{4}{\ohm}} = \SI{3}{\ampere}$

\textbf{Antwortsatz:} Die Stromstärke beträgt $\SI{3}{\ampere}$.

\end{tcolorbox}

\vspace{2em}

\begin{center}
\begin{tikzpicture}
    % Dreieck-Merkregel (URI-Dreieck)
    \draw[thick, physik] (0,0) -- (4,0) -- (2,3) -- cycle;
    % Horizontale Trennlinie (U oben, R|I unten)
    \draw[thick, physik] (0.67,1) -- (3.33,1);
    % Vertikale Trennlinie (R links, I rechts)
    \draw[thick, physik] (2,0) -- (2,1);

    % Beschriftungen
    \node at (2, 1.9) {\Large\textbf{$U$}};
    \node at (1, 0.4) {\Large\textbf{$R$}};
    \node at (3, 0.4) {\Large\textbf{$I$}};

    % Erklärung
    \node[right] at (5, 2) {Zudecken, was gesucht ist:};
    \node[right] at (5, 1.2) {$U$ zudecken $\Rightarrow$ $U = R \cdot I$};
    \node[right] at (5, 0.6) {$R$ zudecken $\Rightarrow$ $R = \frac{U}{I}$};
    \node[right] at (5, 0) {$I$ zudecken $\Rightarrow$ $I = \frac{U}{R}$};

\end{tikzpicture}
\end{center}

% ═══════════════════════════════════════════════════════════════
% ENDE
% ═══════════════════════════════════════════════════════════════

\end{document}
